\begin{frame}
  \onslide<1->{
	\frametitle{Comparison of Process Noise and Measurement Noise} 
	\framesubtitle{$Q$ vs $R$}
}
\begin{itemize}
  \onslide<2->{
\item If $Q$ is large relative to $R$, then $K\approx I$, the extrapolated state estimate will cancel out in the state update 
equation; the updated state estimate will consist essentially only of the measurement $z$
}
  \onslide<3->{
\item This can be seen as follows in the special case that $H(k)$ is square and invertible:
\begin{footnotesize}
\begin{align*}
\stateest{k}{+}&=\stateest{k}{-}+\kalgain{k}\left[\innovation{k}\right]\\
&\approx \stateest{k}{-}+\covmat{k}{-}H(k)\trans\left[H(k)\covmat{k}{-}H(k)\trans\right]^{-1}\left[\innovation{k}\right]\\
&= \stateest{k}{-}+\\
&\;\;\;\;\;\;\;\;\covmat{k}{-}H(k)\trans
\left[H(k)\trans\right]^{-1}\left[\covmat{k}{-}\right]^{-1}H(k)^{-1}\left[\innovation{k}\right]\\
&= \stateest{k}{-}+H(k)^{-1}z(k)-\stateest{k}{-}\\
&= H(k)^{-1}z(k)
\end{align*}
\end{footnotesize}
}
\end{itemize}
\end{frame}


\begin{frame}
  \onslide<1->{
	\frametitle{Comparison of Process Noise and Measurement Noise} 
	\framesubtitle{$Q$ vs $R$}
}
\begin{itemize}
  \onslide<2->{
\item On the other hand, when $R$ dominates $Q$ (and thus $\covmat{k}{-}$), the gain $K$ is nearly 0, 
and the updated estimate is simply the previously extrapolated estimate:
\begin{align*}
\stateest{k}{+}&=\stateest{k}{-}+\kalgain{k}\left[\innovation{k}\right]\\
&\approx \stateest{k}{-}+0\left[\innovation{k}\right]\\
&= \stateest{k}{-}
\end{align*} 
}
  \onslide<3->{
	\item Filter tuning is all about balancing $Q$ and $R$, with $Q$ being less constrained than $R$
}
\end{itemize}
\end{frame}
